% chapter1-beamer.tex — a standalone slide deck for Chapter 1.
%
% This demonstrates philogic's `slim` option.  In slim mode philogic loads ONLY
% the notation (so \nec, \trues, \truthset, \FR, ... are all available), while
% beamer keeps control of theorem blocks, hyperref, and page structure.  The
% same macros you use in the chapter therefore render identically on slides.
%
% Slides are necessarily an *adaptation* of the chapter, not the same source
% \input twice: prose and \section headings do not map onto frames.  Treat this
% as a worked example of turning a chapter into a talk.
%
% Build from this directory (the Makefile target `make chapter1-beamer` does so,
% and sets TEXINPUTS so the bundled philogic.sty in the project root is found).

\documentclass{beamer}

\usetheme{Madrid}
\setbeamertemplate{navigation symbols}{}  % drop the navigation clutter
\setbeamertemplate{theorems}[numbered]

\usepackage[slim]{philogic}

% In slim mode philogic does not load biblatex (a host class usually owns the
% bibliography).  beamer doesn't, so we load it here and reuse the shared
% refs.bib at the project root.
\usepackage[backend=biber]{biblatex}
\addbibresource{../../refs.bib}

\title{Modal Preliminaries}
\subtitle{Chapter 1}
\author{Your Name}
\institute{Your Department \\ University of California, Davis}
\date{}

\begin{document}

\frame{\titlepage}

\begin{frame}{The modal language}
  We work in a modal language \langm{} with propositional variables
  $\mathsf{Prop} = \set{p, q, r, \dots}$, the Boolean connectives, and a unary
  necessity operator $\nec$.  Its dual is
  \[
    \poss A \coloneqq \neg\nec\neg A \qquad (\text{``possibly } A\text{''}).
  \]
\end{frame}

\begin{frame}{Kripke semantics}
  \begin{definition}[Frame]
    A \emph{Kripke frame} is a pair $\FR = \oset{W, R}$ with $W \neq \emptyset$
    and $R \subseteq W \times W$.
  \end{definition}
  \begin{definition}[Model]
    A \emph{Kripke model} $\MRel = \oset{W, R, V}$ adds a valuation
    $V : \mathsf{Prop} \to \power{W}$.  The modal clause:
    \[
      \MRel,w \trues \nec A
      \quad\text{iff}\quad
      \text{for all } v,\ wRv \text{ implies } \MRel,v \trues A.
    \]
  \end{definition}
\end{frame}

\begin{frame}{Normal modal logics}
  The smallest normal modal logic \K{} extends classical propositional logic
  with
  \begin{itemize}
    \item the axiom \kax: $\nec(A \to B) \to (\nec A \to \nec B)$;
    \item the rule \nrule: from $\proves A$, infer $\proves \nec A$.
  \end{itemize}
  Familiar extensions include $\KT = \K \oplus \tax$ (with \tax{} the schema
  $\nec A \to A$) and \Sfour.

  \begin{theorem}[Completeness]
    \K{} is sound and complete with respect to the class of all Kripke frames
    \parencite[ch.~4]{blackburn2001}.
  \end{theorem}
\end{frame}

\begin{frame}[allowframebreaks]{References}
  \printbibliography
\end{frame}

\end{document}
